\begin{frame}{실수 두 번째: ``라벨 없으니 비지도겠지''}
  ``라벨이 없다''는 \textbf{현상}이 ``비지도학습이 맞는 문제''라는
  \textbf{결론}으로 바로 이어지지는 않는다.
  \begin{itemize}
    \item \textbf{``라벨이 비싸서 못 만들었다'' vs ``라벨이
      원리적으로 없다''}: X-ray 폐렴 판별에서 ``폐렴/정상'' 라벨은
      의사가 매겨야 해 비싸고 느리다. 비지도로 푸는 것은 ``라벨을
      만들지 않겠다''는 \emph{선택}이지 ``원리적으로 없다''가 아님
      $\to$ 대개는 ``\textbf{라벨을 일부라도} 만들 수 있는가''를 먼저
      고민 (의사 10장만 라벨 + 나머지 비지도 = \textbf{반지도학습},
      Week 9$\sim$10에서 간접 다룸).
    \item \textbf{``보상이 있긴 한데 '라벨'인 줄 알았다''}:
      ``재구매 여부''는 \textbf{라벨}(지도)로 쓸 수도, ``또 샀으니
      +1''이라는 \textbf{보상}(강화)으로 읽을 수도 --- Q1/Q2로 다시
      확인.
  \end{itemize}
\end{frame}
